\section{Experimental Setup}
\label{sec:experimental-setup}

\subsection{Persisted Benchmark Splits}

The benchmark manifest contains 310 scenarios. Every record stores its scenario
ID, split, grid size, exact blocked cells, start and goal, dynamic objects,
traversal penalties, no-fly configuration, observation-noise level, and
dynamics seed where applicable.

The primary generalization splits each contain 30 scenarios:
\begin{itemize}
    \item the training layout with held-out start/goal pairs;
    \item unseen layouts at obstacle density 0.20;
    \item unseen layouts at density 0.30;
    \item unseen fixed-toggle obstacle locations; and
    \item known toggle locations with unseen periods.
\end{itemize}

Scaling uses ten scenarios each at $15$, $30$, $50$, and $100$ cells per side,
with density 0.20 and new layouts. Realism sweeps independently vary static
density, stochastic toggle probability, moving-obstacle period, spatial energy
penalty, hard versus penalized no-fly zones, and sensor-noise probability.

\subsection{Pairing and Repetitions}

All methods consume the same manifest records. The grid is reconstructed from
stored cells rather than from a method-specific random generator. Classical
planners run on the same graph instance definition; the RL evaluator injects
the identical grid and dynamics into its Gymnasium environment.

Each classical scenario/planner pair and each RL
seed/scenario pair is repeated ten times for timing. Outcome and path metrics
are deterministic under a fixed scenario and policy; repetitions preserve the
raw timing distribution. Route-level compute time is the sum of planner calls
for classical methods and the sum of \texttt{predict()} calls for RL.
Environment construction, stepping, metric calculation, and serialization are
excluded. Per-decision latency is also retained.
Local-observation policies cover all 310 scenarios, producing 15,500 raw rows
per variant ($310\times5\times10$). The fixed-size full-grid ablation is
evaluated only at its trained $15\times15$ input dimension, producing 14,000
rows ($280\times5\times10$); this exclusion is structural rather than a
performance-based omission.

\subsection{Environment and Software Manifest}

The authoritative classical and RL suite manifests both report Linux
6.6.122+, CPython 3.12.13, two logical CPU cores, and no active GPU. Every
manifest stores operating system, CPU, logical core count, Python
implementation and version, package versions, accelerator availability and
names, Git commit, command configuration, manifest hash, and timestamp. The raw
suite was produced in a Drive-backed Google Colab workflow. Millisecond
timings are interpreted only within this recorded runtime class.
The dedicated timing-distribution artifact reports the mean, median, sample
standard deviation, extrema, and repetition count for route-level computation
and per-decision latency, both by policy seed and across seeds.

Historically, one mutable ``latest RL'' environment file was retained rather
than an immutable manifest for every variant; the artifact pipeline now writes
both variant-specific manifests and the compatibility alias. Consequently,
absolute cross-variant timing differences are descriptive, whereas
same-scenario comparisons within a recorded suite are the primary evidence.

\subsection{Statistical Analysis}

Timing repetitions are first collapsed within each policy-seed/scenario pair,
preventing pseudoreplication. RL metrics are then summarized per independent
training seed and reported as the mean and 95\% Student-$t$ confidence
interval across five seeds. Classical confidence intervals use scenarios as
the independent units.

Success differences use exact paired McNemar tests. Path cost and route compute
time use two-sided paired Wilcoxon signed-rank tests with rank-biserial effect
sizes \cite{mcnemar1947note,wilcoxon1945individual}. Only jointly successful
routes enter paired path-cost tests; all routes enter success and compute-time
analyses. Holm correction controls the familywise error rate across the
reported paired tests \cite{holm1979simple}.

\subsection{Reproducibility Checks}

Before evaluation, all persisted scenarios must reconstruct with a valid path
and Dijkstra, A*, and D* Lite must agree on the initial optimal cost. After
evaluation, an integrity pass checks expected row counts, unique run IDs, and
the scenario-manifest hash in every raw result file. Smoke-test checkpoints and
legacy single-seed models are rejected by the expanded evaluator.
